
\begin{figure}[ht]
    \centering
    \begin{minipage}{1.0\textwidth}
        \begin{subfigure}[b]{1.0\textwidth}
            \centering
            \includegraphics[width=1.0\textwidth]{benchmark/figures/video_continues_demo_01.pdf}
            \caption{Text Guidance: A clear acrylic sheet placed on a wooden table with a small dollop of red paint. A rotating paintbrush attached to a rotating platform rotates clockwise and goes through the paint. Static shot with no camera movement.}
            \label{fig:video_continues_demo_01}
        \end{subfigure}
        \begin{subfigure}[b]{1.0\textwidth}
            \centering
            \includegraphics[width=1.0\textwidth]{benchmark/figures/video_continues_demo_02.pdf}
            \caption{Text Guidance: A grabber arm is holding a tennis ball above a piece of cardstock propped up on a rotating platform sitting on a table that rotates clockwise. The grabber lowers the ball and places is on the table as the cardstock rotates. Static shot with no camera movement.
    }
            \label{fig:video_continues_demo_02}
        \end{subfigure}
    \end{minipage}
    \caption{Comparison between video-conditioned (V2V) and image-conditioned (I2V) video continuation. (a) \magi (V2V) accurately captures the pen’s rotational trajectory by leveraging historical motion information, while I2V fails to reproduce the correct motion due to the absence of temporal context. (b) In an occlusion scenario, V2V successfully predicts post-occlusion behavior by utilizing information before the occlusion, whereas I2V shows poor temporal consistency. Each example presents the real-world scene (top row), \magi (V2V) generation (middle row), and \magi (I2V) generation (bottom row).}
\label{fig:video_continues_demo}
\end{figure}
